% !PRINTED 226
\weberpage{226}
Um also endlich die Existenz der Ikosa\"edergruppe festzustellen, ist noch nachzuweisen, dass die Gesammtheit der Substitutionen (20) eine Gruppe bildet. Dies folgt aber aus den nun abzuleitenden Compositionsgesetzen.

Zunächst ergiebt sich sehr einfach aus der Bedeutung von $\Theta$, $\psi$, $\chi$:
\[
\Theta = \begin{pmatrix} \varepsilon, & 0 \\ 0, & 1 \end{pmatrix}, \quad
\psi = \begin{pmatrix} 0, & 1 \\ -1, & 0 \end{pmatrix}, \quad
\chi = \begin{pmatrix} \omega, & 1 \\ 1, & -\omega \end{pmatrix}
\]
\weq{23}{\psi\Theta = \Theta^{-1}\psi, \quad \chi\psi = \psi\chi.}

Ferner erhält man für jeden beliebigen Exponenten $\lambda$:
\[
\chi\Theta^\lambda\chi = \begin{pmatrix} \varepsilon^\lambda\omega^2+1, & (\varepsilon^\lambda-1)\omega \\ (\varepsilon^\lambda-1)\omega, & \varepsilon^\lambda+\omega^2 \end{pmatrix},
\]
oder, indem man nach (5) und (6)
\weq{24}{\omega = \varepsilon+\varepsilon^{-1}, \quad \omega' = \varepsilon^2+\varepsilon^{-2}, \quad \omega\omega' = -1}
setzt:
\[
\chi\Theta^\lambda\chi = \begin{pmatrix} \varepsilon^\lambda\omega-\omega', & \varepsilon^\lambda-1 \\ \varepsilon^\lambda-1, & -\varepsilon^\lambda\omega'+\omega \end{pmatrix}.
\]

Setzen wir darin zunächst $\lambda=1$, so folgt nach (24), (13), (14):
\[
\chi\Theta\chi = \begin{pmatrix} 1-\varepsilon^3, & \varepsilon-1 \\ \varepsilon-1, & \varepsilon-\varepsilon^3 \end{pmatrix}
= \begin{pmatrix} \varepsilon^3(\varepsilon+1), & 1 \\ 1, & -\varepsilon(\varepsilon+1) \end{pmatrix}
\]
\[
= \begin{pmatrix} \varepsilon\omega', & 1 \\ 1, & -\varepsilon^{-1}\omega' \end{pmatrix}
= \begin{pmatrix} -\varepsilon^2, & \varepsilon\omega \\ \varepsilon\omega, & 1 \end{pmatrix}
= \Theta\chi\psi\Theta,
\]
woraus, wenn man beiderseits zur entgegengesetzten Substitution übergeht,
\[
\chi\Theta^{-1}\chi = \Theta^{-1}\chi\psi\Theta^{-1}
\]
folgt. Setzt man andererseits $\lambda=2$, so folgt:
\[
\chi\Theta^2\chi = \begin{pmatrix} \varepsilon-\varepsilon^2, & \varepsilon^2-1 \\ \varepsilon^2-1, & \varepsilon-1 \end{pmatrix}
= \begin{pmatrix} \varepsilon^2(\varepsilon^2+1), & 1 \\ 1, & -\varepsilon(\varepsilon^2+1) \end{pmatrix}
\]
\[
= \begin{pmatrix} \varepsilon^3\omega, & 1 \\ 1, & -\varepsilon^2\omega \end{pmatrix}
= \Theta^{-2}\chi\Theta^{-2},
\]
worin wieder $\Theta$ durch $\Theta^{-1}$ ersetzt werden kann. Man hat daher die Compositionsformeln:

\weq{25}{\chi\Theta^{\pm1}\chi = \Theta^{\pm1}\chi\psi\Theta^{\pm1}, \quad \chi\Theta^{\pm2}\chi = \Theta^{\mp2}\chi\Theta^{\mp2}.}

Durch Anwendung der Formeln (23) und (25) kann man nun je zwei der Substitutionen (20) componiren und gelangt immer wieder auf eine Substitution von der Form (20), wodurch
%CONT>>
